\documentclass[output=paper]{langscibook}
\ChapterDOI{10.5281/zenodo.17641062}
\title{Introduction}
\author{JC Penet\orcid{}\affiliation{Newcastle University, United Kingdom}}
\abstract{\noabstract}
\begin{document}
\maketitle
\section{Why this book?}
Since OpenAI launched ChatGPT back in November 2022, generative artificial intelligence (GenAI) has acted as a major disruptor in the translation industry and beyond. Of course, in any given industry disruptors are not intrinsically good or bad. They merely act as game changers, good and bad. 
Also commonly associated with “chatbots”, GenAI tools certainly are a game changer for the translation industry and, consequently, for translator education too. Using deep learning to train on vast corpora, not only do LLMs have the ability to “generate text which is often indistinguishable from text written by a human” \citep[188]{moorkens_automating_2025}, but they can do much more than just translate texts. In this book, we shall therefore make a distinction between artificial intelligence (AI) and GenAI. In our context, AI refers broadly to all AI technologies, including neural machine translation (NMT). GenAI, however, refers to machine learning tools that generate media, including chat-based Large Language Models (LLMs) like GPT-4. Since this book deals mostly with text, the terms GenAI and LLMs tend to appear synonymously.

Unlike freely available machine translation (MT) tools like Google Translate or DeepL, LLMs can also be prompted to adapt their translations to take the source text and/or the target text context into consideration – for free (initially). In a matter of seconds. It is easy to see how GenAI can potentially contribute to the further democratisation of translation by making it freely accessible in contexts where translation would have been neither practical nor affordable until now. This disruptor also provides translators with new resources that can potentially help them not just to increase the speed and quality of the translation work they deliver to clients, but also with the way they manage their translation projects or go about finding new clients. In the 2024 edition of the annual ELIS (European Language Industry Survey) report, for instance, the independent translators who saw GenAI as a positive trend explained that they leveraged its possibilities “as a tool (e.g. terminology extraction), as a source for editing work, and as a motivator for clients to choose human translation due to bad AI experiences” \citep[24]{elis_elis_2024}.

Of course, the use of AI chatbots for translation work also comes with some important limitations that need our attention. One of them is that they need vast amounts of data to learn, meaning that their ability to translate depends largely on the amount of language data available to train on for any given language. Even though LLMs differ from all other machine learning-based MT systems in that they can be trained on monolingual data instead of bilingual corpora, thus increasing the amount of available data they can train on, they remain mostly trained on English data. Because of this, their translation performance is likely to be much less convincing for so-called low-resource languages compared with high-resource languages . In addition, AI chatbots are trained to identify and replicate language patterns. A natural corollary of this is that they lack “any form of genuine comprehension or the cognitive processes that humans use to understand language and context” \citep[189]{moorkens_automating_2025}. Finally, when GenAI lacks the data it needs to produce part of a translation, it can “hallucinate”, i.e. make up content in the target language. 

These limitations notwithstanding, GenAI has started reshaping what it means to work as a professional translator in an industry that was already largely technology-driven, and where arguably automation has long been used as a way to increase the speed and reach of translation while reducing its cost. This has sometimes been to the detriment of working conditions and job satisfaction for human translators (see, for instance, \cite{lambertwalker2024}). This, in fact, makes it even more urgent for all of us translator educators to engage with GenAI so that we can, in the process, re-examine the role and agency of human translators within the translation process through that lens. In other words, GenAI should encourage us to rethink critically the value and values that humans bring to the translation process. As part of this, we must (re-)interrogate what we do on our translation programmes so that we empower the next generation of professional translators to achieve the kind of “human-centred augmented translation” that will benefit both individuals and society \citep[391]{obrien_human-centered_2024}.

Admittedly, however, for some of us translator educators the prospect of engaging with GenAI may initially have felt somewhat overwhelming. First of all, the timing may not have been ideal. OpenAI’s ChatGPT was released shortly after we came out blinking from a series of pandemic-related lockdowns, during which teaching happened mostly online. For some of us, this meant having to interact with (new) technologies that we may have felt were not always tailored to our own needs and/or that we weren’t always comfortable using. As a result, most of us got to experience first-hand the “technostress” many professional translators experience when they are asked to use translation tools and technologies they are not entirely comfortable with \citep{Penet2024}. This is something we should take seriously, as technostress can “reduc[e] performance and har[m] individual wellbeing” \citep[146]{Koskinen2020}. With its pedagogical approach, this edited volume therefore comes as an attempt to help alleviate feelings of stress among some of us as translator trainers.

Ideally, we should all feel empowered to engage with GenAI on our programmes, where it also acts as a disruptor. Again, this was reflected in the findings of the latest ELIS reports. In its 2024 survey, university staff ranked GenAI implementation as the most widely shared challenge, with close to 90\% of them seeing it as an issue. This prompted the report’s authors to comment that: “Generative AI and how to implement it in the universities' programmes has taken the challenge chart of university staff by storm, topping even their concerns about the visibility of the profession and the eternal lack of time” \citep[27]{elis_elis_2024}. If, the following year, visibility of the profession had reclaimed top spot in the list of challenges, it was still closely followed by GenAI implementation, which remained a concern for over 80\% of university staff \citep[25]{elis_elis_2025}. Yet, adapt and implement we must! This is because, still according to the report, in 2025 “[a]ctual MT use by language companies has increased […] and reaches now the magic mark of 50\% of handled projects. AI makes its entry with an impressive 34\%” \citep[35]{elis_elis_2025}. In a similar vein, Slator, an online portal for language industry news and research, recently made the following point: 
\begin{quote}
    ‘Skate to where the puck is going, not to where it is’ has always been good advice for ice hockey. But it is remarkably prescient for the language services industry as we prepare for the tsunami of disruption brought about by generative artificial intelligence (AI) and large language models, or LLMs. The speed of development has been frenetic since the release of OpenAI’s ChatGPT in November 2022. […] It would be a brave forecaster to predict how the corporate landscape will look five years from now—but operating from where the puck is today risks asking the wrong questions and missing the boat \citep{slator_disruption_2023}.
\end{quote}

Even though this message was clearly intended for language service providers, it certainly holds some relevance for us educators. Whatever we may think or feel about GenAI, we owe it to both our students and to the translation industry to engage with it, and to do so in a way that will help them anticipate where the puck is going. However, we must do so in a way that is ethical and sustainable as, arguably, GenAI is much more than just another tool in the translator’s toolkit. This is what this book intends to help you with.

We have called this edited volume Teaching Translation in the Age of Generative AI: New Paradigm, New Learning. Given the very nature of GenAI, the word paradigm, with its historical meaning of “pattern”, seems rather appropriate. Beyond this, the word captures GenAI’s potential to be a game changer not just in the translation industry, but also in translator education. With this new paradigm, which is characterised by the emergence of new translation practices, comes new learning for both students and educators. This is what we have attempted to capture here with the help of experts from across the globe. In fact, if this edited volume was a film, reviewers would no doubt talk about a “star-studded cast”. As editors, we are truly grateful to all of our colleagues who have so generously accepted the invitation to contribute their expertise, making it available via open access. Like us, they share our ambition to make this book as accessible as possible so it can be of help to as many colleagues in translation education as possible as they learn to negotiate the uncharted territory of GenAI. We are also grateful to Language Science Press and the editors of the “Translation and Multilingual Natural Language Processing” for giving our book a good home. This open access volume builds on outstanding past publications in the series, not least Kenny’s Machine Translation for Everyone: Empowering users in the age of artificial intelligence (\citeyear{kenny2022}). This volume also complements nicely Pym and Hao’s How to Augment Language Skills: Generative AI and Machine Translation in Language Learning and Translator Training (\citeyear{pym_how_2025}) as it focuses more specifically on GenAI in translator education. 

\section{Contents of this book}
This book is articulated around three main parts. Part 1 explores the new skills and competences we need to help our students develop in the age of GenAI. It opens with Gary Massey and Maureen Ehrensberger-Dow’s detailed discussion in Chapter 1 of translation competence in the age of GenAI. They make the case that GenAI is a collaborative technology that can be complemented by human agency if human translators bring their own, complementary skills to the table. To that aim, translation programmes should foreground textual, digital skills and interlingual skills when training students to work with GenAI. In Chapter 2, Erik Angelone discusses the importance of deliberate practice in translator training to help students develop the competence and expertise they will need in their future careers. He then shows how immediate, informative feedback is a key condition for the development of deliberate practice and how we can use AI chatbots in our teaching as More Knowledgeable Others (MKO) that provide students with the required feedback to foster translational expertise through deliberate practice. Chapter 3, by Ramón Inglada, argues that, in the age of automated translation, the concept suitability should take precedence over that of correctness. To make sure our students are in a position to evaluate the suitability of the translation solutions generated by (Gen)AI, he argues, we should help them develop three new skills, namely the skills of selection, assessment and prompting, alongside the traditional core skills for translation competence. He then gives concrete examples of how GenAI can be used in the classroom to foster these new skills among trainee translators.

In Part 2, the focus shifts to the new knowledge we and our students need to delve into in the age of GenAI. This knowledge often underpins the new skills and competences mentioned in Part 1. In Chapter 4, for instance, Lynne Bowker argues the importance of teaching our translation students about data (and, therefore, corpora) as it can help them make more informed and responsible use of the technology. She recognises, however, that doing so from an expert-to-expert position may not always be most effective, as it is easy to get lost in a rabbit-hole of intricate details when talking about AI networks. Instead, she shows through concrete and compelling examples that borrowing techniques from science communication – whose goal is to make expert communication available to non-experts – can make our teaching about data much more impactful for students. This is followed by Chapter 5, in which Masaru Yamada discusses prompting for GenAI. More specifically, he shows through detailed examples the usefulness of traditional translation studies concepts to prompt AI chatbots, thus opening new avenues for translator education. In an echo to Bowker’s chapter, Joss Moorkens and Gökhan Doğru tackle the all-important topic of AI ethics in Chapter 6. They consider ways to introduce discussion and reflection in the classroom that maximise the impact of this teaching. To that end, they focus in particular on case studies and the use of mapping and images to foster systemic thinking around translation and GenAI.

Finally, Part 3 puts some flesh on the bones of Parts 1 and 2 as it reviews some of the new teaching approaches adopted by colleagues since the advent of GenAI. It does so by introducing the reader to a series of vignettes taken from a variety of translation-related disciplines and contexts. Chapter 7 opens with Senem Öner Bulut presenting the results of a research project in which staff and students collaboratively and experimentally explored the dynamics of prompt engineering for translation, as well as the new skills human translators may need to develop to prompt effectively. In Chapter 8, Maria Zimina-Poirot gives concrete examples of activities to develop students’ critical and ethical use of GenAI as part of the translation workflow. Through these, she convincingly shows how such a use of GenAI can help augment human contributions. This is followed by Chapter 9, in which David Orrego-Carmona moves the lens to subtitler training as another area of translator education where “the teaching must shift from teaching specific technical and translation skills to developing adaptable professionals who can critically engage with evolving technologies while maintaining high-quality standards”. To that aim, he presents us with series of innovative classroom activities that can be replicated to help our students’ “digital reflexivity” in subtitling. In Chapter 10, Sonia Vandepitte focuses on another specific area of translator education, namely translation into one’s second language (L2 translation). Reporting on the findings of an exploratory research project involving students of an L2 translation class on a postgraduate programme, she shows how integrating GenAI in L2 translation at the editing stage of the process not only enhances students’ reflexivity, but it can also increase their intrinsic motivation, which is crucial in L2 translation training. In Chapter 11, we move to interpreting as Sahar Othomani and Nermin al Sharman introduce us to the [AI]phra project, a simulation software created to enhance interpreter training in the MENA (Middle East and North Africa) region through gamification. Piloted in 2023, the key motivator behind [AI]phra is to unleash the potential of AI to address the at-times limited access to interpreter education in the MENA region. Finally, in Chapter 12 we move slightly away from translator education to explore the potential of embracing MT in second language (L2) education. In this final chapter, Atsushi Mizumoto thus introduces us to the Metacognitive Resource Use Framework, which “positions learners as metacognitive agents capable of strategically utilizing a wide range of language resources, including MT and generative AI tools”.

\section{One final thought}
As GenAI continues to reshape what it means to translate and, therefore, what it means to be a translator (educator), our main aim with this volume is to help fellow translator educators better understand the new skills and competences that we need to foster in our teaching, scaffolded by new knowledge. We also hope that, taking a leaf from some of the case studies in Part 3, many colleagues will be spurred on to further explore what it means to (learn to) translate in the age of GenAI in partnership with students. This, in turn, could lead to a welcome increase in the number of research publications on the topic at a time where we all need to learn to navigate our way through this new paradigm.


{\sloppy\printbibliography[heading=subbibliography]}
\end{document}

